\chapter{Mittelschmerz (Ovulation Pain)}

\section*{Definition}
Mittelschmerz is mid-cycle abdominal pain that occurs at the time of ovulation, approximately 14 days before the next menstrual period. It is a normal physiologic phenomenon that is typically mild and self-limiting.

\section*{What Happens in Your Body}
At ovulation, the mature ovarian follicle ruptures to release the egg, accompanied by a small amount of blood and follicular fluid. This fluid irritates the peritoneum, causing localized lower abdominal pain. The pain typically alternates sides each month depending on which ovary ovulates.

\section*{Causes}
\begin{itemize}
\item Normal ovulation (physiologic)
\item Ovarian follicle rupture releasing fluid and blood
\item Peritoneal irritation from follicular contents
\end{itemize}

\section*{When to See a Doctor}
\begin{itemize}
\item Lower abdominal pain midway through the menstrual cycle
\item Pain typically on one side (alternating each cycle)
\item Dull, cramping, or sharp pain lasting minutes to hours
\item May be accompanied by light spotting
\item No other symptoms of infection or pathology
\item Pain resolves within 24-48 hours
\end{itemize}

\section*{Related Symptoms}
\begin{itemize}
\item Menstrual cramps (dysmenorrhea)
\item Ovarian cyst rupture (more severe pain)
\item Appendicitis (right-sided pain, must be distinguished)
\item Ectopic pregnancy (must be ruled out if sexually active)
\end{itemize}

\section*{Self-Care / Home Management}
\begin{itemize}
\item Apply a heating pad to the lower abdomen
\item Take OTC pain relievers (ibuprofen, acetaminophen)
\item Rest until the pain subsides
\item Track ovulation and mittelschmerz on a calendar
\item Pain relief with warm baths
\item Practice relaxation breathing
\end{itemize}

\section*{How Your Doctor Will Evaluate You}
\begin{itemize}
\item Clinical diagnosis based on timing and character of pain
\item Consistent relationship to menstrual cycle
\item Pelvic examination is normal
\item Ultrasound if pain is severe or atypical to rule out ovarian cyst or other pathology
\item Urine pregnancy test to rule out ectopic pregnancy
\end{itemize}

\section*{Treatment Options}
\begin{itemize}
\item Reassurance that mittelschmerz is a normal, benign condition
\item NSAIDs or acetaminophen for pain
\item Combined oral contraceptives suppress ovulation and eliminate mittelschmerz
\item Heating pad and rest
\end{itemize}

\section*{References}
\begin{thebibliography}{99}
\bibitem{mittelschmerz:ref1} Lippmann S. ``Mittelschmerz.'' \textit{American Family Physician}, 1995;51(6):1541--1542.
\bibitem{mittelschmerz:ref2} Berek JS, Novak E. \textit{Berek \& Novak's Gynecology}, 16th edition. Wolters Kluwer, 2019.
\bibitem{mittelschmerz:ref3} Deligeoroglou E, Creatsas G. ``Menstrual disorders.'' \textit{Annals of the New York Academy of Sciences}, 2000;900(1):247--254.
\bibitem{mittelschmerz:ref4} Crottogini JJ, Finkielman S. ``Mittelschmerz: a study of 100 cases.'' \textit{Obstetrics \& Gynecology}, 1967;29(6):777--781.
\bibitem{mittelschmerz:ref5} Hatcher RA, Trussell J, Nelson AL, et al. \textit{Contraceptive Technology}, 21st edition. Ayer Company Publishers, 2018.
\bibitem{mittelschmerz:ref6} Fritz MA, Speroff L. \textit{Clinical Gynecologic Endocrinology and Infertility}, 8th edition. Wolters Kluwer, 2011.

\end{thebibliography}
